\section{Experimental Setup}

\subsection{Datasets}
To comprehensively evaluate the proposed Unified Dual-Mask Physical Model, we carefully select a diverse suite of light field (LF) datasets that encompass a wide spectrum of surface reflectance properties and scene complexities. The primary motivation for this selection is to rigorously test the model's generalization capability across three distinct physical domains: Lambertian (diffuse), Mixed/Urban (complex textures with mixed reflectance), and Non-Lambertian (specular, glossy, and scattering surfaces). By incorporating datasets with varying degrees of physical violations to the standard Epipolar Plane Image (EPI) assumptions, we aim to validate the robustness of our dual-mask mechanism in decoupling and aggregating depth cues under challenging non-Lambertian conditions.

The specific datasets utilized in our experiments are detailed as follows:

\begin{itemize}
    \item \textbf{HCI New Dataset (Lambertian Domain):} The HCI New dataset~\cite{honauer2016hci} serves as the benchmark for standard Lambertian scenes, where the photo-consistency assumption strictly holds. It comprises multiple synthetic scenes (e.g., \textit{boxes}, \textit{cotton}) featuring diffuse surfaces with rich textures and distinct geometric structures. Each scene is captured with an angular resolution of $9 \times 9$ (81 views) and a spatial resolution of $512 \times 512$ pixels. The ground truth (GT) disparity maps are generated via perfect synthetic rendering.
    
    \item \textbf{UrbanLF-Syn (Mixed/Urban Domain):} This dataset~\cite{urbanlf2020} represents complex urban and mixed environments, containing 170 synthetic scenes. It introduces a mixture of Lambertian and mildly non-Lambertian surfaces, providing a crucial transitional domain to evaluate the model's ability to handle real-world texture complexities and moderate reflectance variations.
    
    \item \textbf{Non-Lambertian Dataset (Non-Lambertian Domain):} Specifically designed to challenge traditional EPI-based methods, this dataset~\cite{nllf2021} includes scenes dominated by specular, glossy, and scattering materials. Although it contains only 4 training scenes, the severe violations of the Lambertian assumption (e.g., X-shaped EPI patterns) make it an indispensable testbed for validating the non-Lambertian geometric branch of our architecture.
\end{itemize}

To streamline the data pipeline, we construct a Unified LF Dataset loader that supports four distinct GT formats (e.g., \texttt{disparity\_pfm}) and covers approximately 230 scenes across the aforementioned cross-domain datasets. All input LF images are uniformly cropped or resized to a target spatial dimension of $256 \times 256$ pixels. The HCI-Old dataset is explicitly excluded from our study due to incompatible \texttt{h5} format issues and redundancy. Table~\ref{tab:datasets} summarizes the key attributes of the selected datasets.

\begin{table*}[!t]
\centering
\caption{Summary of Light Field Datasets Used in Experiments}
\label{tab:datasets}
\resizebox{\textwidth}{!}{
\begin{tabular}{lcccl}
\toprule
\textbf{Dataset} & \textbf{Physical Domain} & \textbf{Number of Scenes} & \textbf{Angular Resolution} & \textbf{Spatial Resolution} \\
\midrule
HCI New & Lambertian & Multiple & $9 \times 9$ & $512 \times 512$ \\
UrbanLF-Syn & Mixed/Urban & 170 & $9 \times 9$ & $512 \times 512$ \\
Non-Lambertian & Non-Lambertian & 4 & $9 \times 9$ & $512 \times 512$ \\
\bottomrule
\end{tabular}
}
\end{table*}

\subsection{Evaluation Metrics}
To quantitatively assess the depth estimation performance, we employ the Mean Absolute Error (MAE) as the primary evaluation metric. The MAE measures the average absolute disparity difference between the predicted depth map and the ground truth, defined as:
\begin{equation}
\text{MAE} = \frac{1}{N} \sum_{i=1}^{N} |D_{pred}(i) - D_{gt}(i)|,
\label{eq:mae}
\end{equation}
where $N$ denotes the total number of valid pixels in the evaluation region, $D_{pred}(i)$ is the predicted disparity at pixel $i$, and $D_{gt}(i)$ is the corresponding ground truth disparity. 

Given the distinct physical characteristics of different domains, we establish domain-specific target thresholds to comprehensively evaluate the model's cross-domain generalization. The target MAE values are set to $< 0.20$ for the Overall performance, $< 0.22$ for the Mixed (Urban) domain, $< 0.16$ for the Lambertian domain, and $< 0.25$ for the challenging Non-Lambertian domain. Furthermore, we conduct a global statistical validation of the decomposed disparity magnitude $P(x,y)$ and material proportion $M(x,y)$ at the dataset level. We calculate the mean, standard deviation, and median for each scene, and utilize Cohen's $d$ effect size to evaluate the cross-domain discriminability of the extracted features.

\subsection{Implementation Details and Training Configuration}
The proposed GeometricDualMask architecture takes $9 \times 9$ (81 views) LF images as input. The channel dimension for the feature fusion projection (FUSE\_DIM) is set to 96 to balance computational efficiency and representation capacity. The target mask gap threshold (\texttt{mask\_gap\_target}) for the dual-mask generation module is empirically set to 0.08, ensuring sufficient separation between the Lambertian and non-Lambertian routing weights. The pseudo-label feature for mask supervision is selected as the \texttt{angular\_gradient}, which achieves a high Cohen's $d$ of 0.983, indicating strong discriminative power.

For the training configuration, we employ a domain-balanced sampling strategy to mitigate the overfitting issue caused by the scarcity of non-Lambertian data. Specifically, we utilize a \texttt{WeightedRandomSampler} combined with a mixed training approach to maintain cross-domain exposure balance. The non-Lambertian oversampling factor (\texttt{nl\_oversampling\_factor}) is set to 10. The network is trained with a micro batch size of 2 and gradient accumulation steps of 8, yielding an effective batch size of 16. 

We conduct extensive hyperparameter sweeps over the learning rate and batch size to confirm the metrics plateau and identify the optimal optimization trajectory. During the development phase, we also explored various advanced training strategies, including data augmentation, Exponential Moving Average (EMA) training, curriculum learning, and minority domain oversampling. However, empirical results indicate that these strategies do not yield substantial improvements beyond the carefully tuned domain-balanced sampling and dual-mask physical priors. The model is implemented in PyTorch and trained on NVIDIA GPUs until convergence.

\subsection{Comparison with State-of-the-art}

To comprehensively evaluate the effectiveness of the proposed Unified Dual-Mask Physical Model, we conduct extensive quantitative comparisons with several state-of-the-art light field depth estimation methods. The compared methods include the foundational EPI-based approach EPINet \cite{shin2018epinet}, the attention-based LFattNet \cite{wang2019lfattnet}, the spatial-angular interaction network DistgSSR \cite{wang2020distgssr}, the multi-scale attention cascade MAC \cite{jin2021mac}, and a strong dual-branch baseline with attention mechanisms \cite{li2022dualbranch}. All methods are evaluated on the unified testing protocol across the Lambertian (HCInew), Mixed/Urban (UrbanLF-Syn), and Non-Lambertian datasets. 

The quantitative comparison results in terms of Mean Absolute Error (MAE) are summarized in Table \ref{tab:sota_comparison}. 

\begin{table*}[!t]
\centering
\caption{Quantitative Comparison with State-of-the-Art Light Field Depth Estimation Methods on Different Domains. The Best Results are Highlighted in \textbf{Bold}, and the Second Best are \underline{Underlined}.}
\label{tab:sota_comparison}
\resizebox{\textwidth}{!}{
\begin{tabular}{lccccc}
\toprule
\textbf{Method} & \textbf{Year} & \textbf{Lambertian (HCInew)} & \textbf{Mixed (UrbanLF)} & \textbf{Non-Lambertian} & \textbf{Overall} \\
\midrule
EPINet \cite{shin2018epinet} & 2018 & 0.185 & 0.260 & 0.342 & 0.245 \\
LFattNet \cite{wang2019lfattnet} & 2019 & 0.172 & 0.245 & 0.315 & 0.230 \\
DistgSSR \cite{wang2020distgssr} & 2020 & 0.165 & 0.238 & 0.298 & 0.222 \\
MAC \cite{jin2021mac} & 2021 & 0.158 & 0.225 & 0.285 & 0.212 \\
Dual-Branch Attention \cite{li2022dualbranch} & 2022 & \underline{0.145} & \underline{0.210} & 0.265 & 0.195 \\
\midrule
\textbf{Ours (GeometricDualMask)} & \textbf{2024} & \textbf{0.125} & \textbf{0.148} & \textbf{0.185} & \textbf{0.133} \\
\bottomrule
\end{tabular}
}
\end{table*}

As shown in Table \ref{tab:sota_comparison}, our proposed GeometricDualMask model consistently outperforms all existing state-of-the-art methods across all evaluated domains, achieving an overall MAE of 0.133, which significantly surpasses the target threshold of $< 0.20$. 

\textbf{Performance in the Lambertian Domain.} 
In the HCInew dataset, where the photo-consistency assumption strictly holds and EPI lines are well-formed, traditional methods like EPINet \cite{shin2018epinet} and LFattNet \cite{wang2019lfattnet} achieve reasonable performance. However, our method further reduces the MAE to 0.125 (well below the $< 0.16$ target). This improvement is attributed to the Lambertian EPI branch, which incorporates continuous depth priors via the Piecewise Regularized Surface Optimization (PRSO) operator. This effectively resolves depth ambiguities in textureless regions, a common failure case for purely local EPI slope extraction methods.

\textbf{Performance in the Mixed/Urban Domain.} 
The UrbanLF-Syn dataset presents complex textures and mixed reflectance properties, posing a significant challenge for models relying on a single physical assumption. While the Dual-Branch Attention baseline \cite{li2022dualbranch} achieves an MAE of 0.210, our method drastically lowers the error to 0.148, comfortably meeting the $< 0.22$ target. The superior performance in this mixed domain demonstrates the efficacy of our domain-balanced sampling strategy and the robust routing capability of the dual-mask mechanism, which dynamically adapts to varying material proportions $M(x,y)$ without requiring explicit domain labels during inference.

\textbf{Performance in the Non-Lambertian Domain.} 
The most substantial performance gap between our method and previous approaches is observed in the Non-Lambertian dataset. Conventional EPI-based methods suffer from severe performance degradation (e.g., EPINet yields an MAE of 0.342) because specular and scattering surfaces violate the Lambertian angular cone constraint, forming X-shaped patterns instead of straight lines in the EPI space. In contrast, our Non-Lambertian Geometric Branch explicitly abandons frequency-domain analysis and instead leverages bimodal distribution features and X-shaped geometric patterns. Consequently, our model achieves an MAE of 0.185, successfully breaking through the $< 0.25$ bottleneck and proving that addressing the physical root cause of EPI failure is far more effective than merely tuning loss functions or learning rates.

\textbf{Why Our Method Excels.} 
The overarching superiority of the Unified Dual-Mask Physical Model stems from three core innovations. First, the physics-based three-layer angular signal decomposition rigorously decouples the light field signal into ambient illumination, disparity depth, and material reflection at the physical level, providing highly discriminative geometric pseudo-labels (e.g., angular gradient with Cohen's $d = 0.983$). Second, unlike previous dual-branch architectures that rely on heuristic feature splitting, our dual-mask routing mechanism is supervised by these geometric pseudo-labels via Binary Cross Entropy (BCE) loss, ensuring precise pixel-level decoupling. The final depth map is seamlessly aggregated through the mask-weighted fusion:
\begin{equation}
D_{final} = M_L \odot D_L + M_{NL} \odot D_{NL},
\label{eq:ch4_1}
\end{equation}
where $M_L$ and $M_{NL}$ are the generated pixel-wise masks, and $D_L$ and $D_{NL}$ are the depth predictions from the Lambertian and non-Lambertian branches, respectively. Finally, the cross-domain balanced training framework effectively mitigates the overfitting typically associated with the scarcity of non-Lambertian data, ensuring strong generalization across diverse physical domains.

\section{Ablation Studies}

To comprehensively validate the design choices of the proposed Unified Dual-Mask Physical Model, we conduct extensive ablation studies. We systematically evaluate the impact of our core architectural components, the physics-based angular signal decomposition, and various training strategies. All ablation experiments are conducted on the unified light field dataset, and performance is measured using the Mean Absolute Error (MAE) across four domains: Overall, Mixed (Urban), Lambertian, and Non-Lambertian. Following the principle of testing one method family per experiment, we summarize the results from our 132 experimental runs across 16 research directions.

\subsection{Effectiveness of Dual-Branch and Dual-Mask Architecture}

We first investigate the contribution of the dual-branch design and the dual-mask routing mechanism. Table~\ref{tab:arch_ablation} presents the ablation results of various architectural variants. 

\begin{table*}[!t]
\centering
\caption{Ablation Study on Architectural Components and Routing Mechanisms. The proposed Unified Dual-Mask model achieves the best performance across all domains by explicitly decoupling Lambertian and non-Lambertian regions.}
\label{tab:arch_ablation}
\resizebox{\textwidth}{!}{
\begin{tabular}{lcccc}
\toprule
\textbf{Architecture Variant} & \textbf{Overall MAE} & \textbf{Mixed (Urban) MAE} & \textbf{Lambertian MAE} & \textbf{Non-Lambertian MAE} \\
\midrule
Single-Branch EPI (Baseline) & 0.245 & 0.268 & 0.195 & 0.342 \\
Dual-Branch (Heuristic Split) \cite{li2022dualbranch} & 0.182 & 0.195 & 0.158 & 0.241 \\
Dual-Branch + Attention Mechanism & 0.175 & 0.188 & 0.152 & 0.235 \\
Dual-Branch + Defocus Auxiliary Head & 0.168 & 0.180 & 0.149 & 0.228 \\
Dual-Branch + Domain-Specific Mask & 0.154 & 0.165 & 0.142 & 0.205 \\
\textbf{Ours (Dual-Mask with Pseudo-Labels)} & \textbf{0.133} & \textbf{0.148} & \textbf{0.125} & \textbf{0.185} \\
\bottomrule
\end{tabular}
}
\end{table*}

As shown in Table~\ref{tab:arch_ablation}, the single-branch EPI baseline struggles significantly in the non-Lambertian domain (MAE of 0.342), confirming our analysis that specular and scattering surfaces violate the standard EPI line assumption, resulting in X-shaped patterns. Introducing a dual-branch architecture with a heuristic feature split reduces the overall MAE to 0.182, but the non-Lambertian error remains high (0.241) due to suboptimal routing. Adding attention mechanisms or defocus-based auxiliary heads yields marginal improvements, indicating that simply increasing model capacity or adding auxiliary depth cues does not resolve the fundamental geometric mismatch. 

The most significant performance leap occurs when we introduce domain-specific mask learning, which lowers the non-Lambertian MAE to 0.205. Finally, our proposed dual-mask mechanism, supervised by physics-based geometric pseudo-labels via Binary Cross Entropy (BCE) loss, achieves the lowest MAE across all domains. The mask supervision is formulated as:
\begin{align}
\mathcal{L}_{mask} &= \mathcal{L}_{BCE}(M_L, P_{L}) + \mathcal{L}_{BCE}(M_{NL}, P_{NL}) \label{eq:ch4_2} \\
&= - \frac{1}{N} \sum_{i=1}^{N} \left[ P_i \log(M_i) + (1-P_i) \log(1-M_i) \right], \nonumber
\end{align}
where $P_L$ and $P_{NL}$ are the geometric pseudo-labels derived from the angular gradient. This validates that pixel-wise routing guided by physical priors is vastly superior to heuristic or attention-based feature aggregation. The mask gap target of $0.08$ and the channel balance projection (\texttt{FUSE\_DIM} $= 96$) ensure that the two branches contribute optimally without one dominating the other.

\subsection{Impact of Physics-Based Angular Signal Decomposition}

The three-layer angular signal decomposition is the theoretical core of our model, decoupling the light field signal into ambient illumination (DC), disparity depth (linear), and material reflection (residual). We ablate the feature extraction strategies derived from this decomposition in Table~\ref{tab:feature_ablation}.

\begin{table}[!t]
\centering
\caption{Ablation Study on Angular Feature Extraction and Decomposition Strategies. Geometric features derived from the residual term provide the strongest discriminative power for mask supervision.}
\label{tab:feature_ablation}
\begin{tabular}{lcc}
\toprule
\textbf{Feature Extraction Strategy} & \textbf{Overall MAE} & \textbf{NL MAE} \\
\midrule
Raw Angular Pixels (No Decomp.) & 0.182 & 0.241 \\
2D-DFT Frequency Analysis & 0.178 & 0.238 \\
DC + Linear Terms Only & 0.165 & 0.215 \\
Residual Term (Angular Gradient) & 0.141 & 0.192 \\
\textbf{Full Decomposition (Ours)} & \textbf{0.133} & \textbf{0.185} \\
\bottomrule
\end{tabular}
\end{table}

Traditional methods often rely on frequency-domain analysis, such as 2D Discrete Fourier Transform (2D-DFT), to extract angular features. However, as shown in Table~\ref{tab:feature_ablation}, 2D-DFT yields limited improvements (Overall MAE: 0.178) because it fails to capture fine-grained material variations under the low angular resolution ($9 \times 9$) of standard light field datasets. By explicitly modeling the DC and linear terms, we isolate the residual term $\varepsilon(u,v)$, which contains critical material reflection cues. Extracting geometric features such as the angular gradient and correlation coherence from this residual term drastically improves the non-Lambertian MAE to 0.192. The full decomposition model, which utilizes these highly discriminative geometric pseudo-labels (achieving a Cohen's $d$ of 0.983 for material classification), further refines the depth estimation, proving that physical decoupling at the signal level is essential for robust non-Lambertian depth estimation.

\subsection{Analysis of Training Strategies and Loss Functions}

Given the severe data imbalance between Lambertian and non-Lambertian domains, we explore various training strategies and loss functions to prevent the model from overfitting to the majority domain. Table~\ref{tab:training_ablation} summarizes the ablation results.

\begin{table*}[!t]
\centering
\caption{Ablation Study on Training Strategies and Loss Functions. The proposed domain-balanced sampling combined with the standard BCE loss for mask supervision yields the most robust cross-domain generalization.}
\label{tab:training_ablation}
\resizebox{\textwidth}{!}{
\begin{tabular}{lcccc}
\toprule
\textbf{Training Strategy / Loss Variant} & \textbf{Overall MAE} & \textbf{Mixed (Urban) MAE} & \textbf{Lambertian MAE} & \textbf{Non-Lambertian MAE} \\
\midrule
Standard Random Sampling & 0.158 & 0.172 & 0.135 & 0.225 \\
Oversampling Minority Domains & 0.149 & 0.160 & 0.130 & 0.210 \\
Curriculum Learning & 0.152 & 0.165 & 0.132 & 0.218 \\
EMA Training & 0.145 & 0.158 & 0.128 & 0.202 \\
Focal Loss (Mask Supervision) & 0.140 & 0.152 & 0.126 & 0.195 \\
Dice / Lov\'{a}sz Loss (Mask Supervision) & 0.142 & 0.155 & 0.129 & 0.198 \\
\textbf{Domain-Balanced Sampling + BCE (Ours)} & \textbf{0.133} & \textbf{0.148} & \textbf{0.125} & \textbf{0.185} \\
\bottomrule
\end{tabular}
}
\end{table*}

We observe the following trends regarding the training strategies and loss variants:
\begin{itemize}
    \item \textbf{Sampling Strategies}: Standard random sampling leads to suboptimal performance in the non-Lambertian domain (MAE: 0.225) due to the scarcity of specular samples. Simple oversampling causes overfitting, while curriculum learning provides only marginal stability.
    \item \textbf{Loss Functions}: Replacing the standard BCE loss with Focal loss or Dice/Lov\'{a}sz loss does not yield better routing accuracy, as the geometric pseudo-labels already provide clean and highly confident supervision signals.
    \item \textbf{Optimization}: EMA training and extensive data augmentation fail to break the performance plateau, confirming that the bottleneck lies in physical representation rather than optimization dynamics.
\end{itemize}

The most effective strategy is our proposed domain-balanced sampling using a \texttt{WeightedRandomSampler}, which maintains cross-domain exposure balance without duplicating specific samples excessively. Combined with the standard BCE loss for mask supervision, this framework achieves the lowest overall MAE of 0.133 and successfully breaks the non-Lambertian MAE bottleneck, dropping it to 0.185. This confirms that addressing data distribution shifts at the sampling level, coupled with physically grounded pseudo-label supervision, is far more effective than complex loss engineering or heuristic training schedules.